
\documentclass{article}
\usepackage{colortbl}
\usepackage{makecell}
\usepackage{multirow}
\usepackage{supertabular}

\begin{document}

\newcounter{utterance}

\centering \large Interaction Transcript for game `bbh', experiment `logical\_deduction\_five\_objects', episode 91 with selfrefine{-}Qwen3.5{-}27B.
\vspace{24pt}

{ \footnotesize  \setcounter{utterance}{1}
\setlength{\tabcolsep}{0pt}
\begin{supertabular}{c@{$\;$}|p{.15\linewidth}@{}p{.15\linewidth}p{.15\linewidth}p{.15\linewidth}p{.15\linewidth}p{.15\linewidth}}
    \# & \multicolumn{2}{c}{Player} && \multicolumn{2}{c}{Game Master} \\
    \hline

    \theutterance \stepcounter{utterance}  
    & & & \multicolumn{4}{p{0.6\linewidth}}{
        \cellcolor[rgb]{0.9,0.9,0.9}{
            \makecell[{{p{\linewidth}}}]{
                \texttt{\tiny{[P1$\langle$GM]}}
                \texttt{A logical deduction task which requires deducing the order of a sequence of objects.} \\
\\ 
\texttt{Q: The following paragraphs each describe a set of three objects arranged in a fixed order. The statements are logically consistent within each paragraph. In a golf tournament, there were three golfers: Amy, Eli, and Eve. Eve finished above Amy. Eli finished below Amy.} \\
\texttt{Options:} \\
\texttt{(A) Amy finished last} \\
\texttt{(B) Eli finished last} \\
\texttt{(C) Eve finished last} \\
\texttt{A: (B)} \\
\\ 
\texttt{Q: The following paragraphs each describe a set of three objects arranged in a fixed order. The statements are logically consistent within each paragraph. On a shelf, there are three books: a white book, a green book, and an orange book. The green book is to the right of the white book. The orange book is the rightmost.} \\
\texttt{Options:} \\
\texttt{(A) The white book is the leftmost} \\
\texttt{(B) The green book is the leftmost} \\
\texttt{(C) The orange book is the leftmost} \\
\texttt{A: (A)} \\
\\ 
\texttt{Q: The following paragraphs each describe a set of three objects arranged in a fixed order. The statements are logically consistent within each paragraph. On a shelf, there are three books: a red book, a gray book, and a white book. The white book is to the left of the gray book. The red book is the second from the left.} \\
\texttt{Options:} \\
\texttt{(A) The red book is the leftmost} \\
\texttt{(B) The gray book is the leftmost} \\
\texttt{(C) The white book is the leftmost} \\
\texttt{A: (C)} \\
\\ 
\texttt{Q: The following paragraphs each describe a set of five objects arranged in a fixed order. The statements are logically consistent within each paragraph. On a branch, there are five birds: a cardinal, a robin, a blue jay, a quail, and a raven. The robin is to the right of the raven. The cardinal is the leftmost. The raven is to the right of the blue jay. The blue jay is the third from the left.} \\
\texttt{Options:} \\
\texttt{(A) The cardinal is the rightmost} \\
\texttt{(B) The robin is the rightmost} \\
\texttt{(C) The blue jay is the rightmost} \\
\texttt{(D) The quail is the rightmost} \\
\texttt{(E) The raven is the rightmost} \\
\texttt{A:} \\
            }
        }
    }
     \\ \\

    \theutterance \stepcounter{utterance}  
    & & \multicolumn{4}{p{0.6\linewidth}}{
        \cellcolor[rgb]{0.95,0.95,0.95}{
            \makecell[{{p{\linewidth}}}]{
                \texttt{\tiny{[P1$\leftrightarrow$P]}}
                \texttt{{[}DRAFT{]}} \\
\texttt{To determine the correct answer, let's break down the positions of the five birds based on the clues provided. We have five positions, numbered 1 (leftmost) to 5 (rightmost).} \\
\\ 
\texttt{1.  **"The cardinal is the leftmost."**} \\
\texttt{    *   Position 1: Cardinal} \\
\\ 
\texttt{2.  **"The blue jay is the third from the left."**} \\
\texttt{    *   Position 3: Blue Jay} \\
\\ 
\texttt{3.  **"The raven is to the right of the blue jay."**} \\
\texttt{    *   Since the Blue Jay is at position 3, the Raven must be at position 4 or 5.} \\
\\ 
\texttt{4.  **"The robin is to the right of the raven."**} \\
\texttt{    *   This implies the order is: Blue Jay < Raven < Robin.} \\
\texttt{    *   Since the Blue Jay is at position 3, the Raven must be at position 4 (to allow space for the Robin to its right).} \\
\texttt{    *   Consequently, the Robin must be at position 5.} \\
\texttt{    *   Current arrangement:} \\
\texttt{        *   1: Cardinal} \\
\texttt{        *   2: ?} \\
\texttt{        *   3: Blue Jay} \\
\texttt{        *   4: Raven} \\
\texttt{        *   5: Robin} \\
\\ 
\texttt{5.  **Determine the remaining bird:**} \\
\texttt{    *   The only bird left is the **quail**, and the only empty spot is position 2.} \\
\texttt{    *   Final arrangement:} \\
\texttt{        1.  Cardinal} \\
\texttt{        2.  Quail} \\
\texttt{        3.  Blue Jay} \\
\texttt{        4.  Raven} \\
\texttt{        5.  Robin} \\
\\ 
\texttt{The question asks to identify the rightmost bird. Based on the deduction, the bird at position 5 is the **robin**.} \\
\\ 
\texttt{Comparing this to the options:} \\
\texttt{(A) The cardinal is the rightmost (False, it is leftmost)} \\
\texttt{(B) The robin is the rightmost (True)} \\
\texttt{(C) The blue jay is the rightmost (False, it is third)} \\
\texttt{(D) The quail is the rightmost (False, it is second)} \\
\texttt{(E) The raven is the rightmost (False, it is fourth)} \\
\\ 
\texttt{A: (B)} \\
            }
        }
    }
    & \\ \\

    \theutterance \stepcounter{utterance}  
    & & \multicolumn{4}{p{0.6\linewidth}}{
        \cellcolor[rgb]{0.95,0.95,0.95}{
            \makecell[{{p{\linewidth}}}]{
                \texttt{\tiny{[P1$\leftrightarrow$P]}}
                \texttt{{[}SELF{-}FEEDBACK{]}} \\
\texttt{FORMAT: problem} \\
\texttt{RULES: ok} \\
\texttt{STRATEGY: ok} \\
\texttt{REASON: The answer draft includes a detailed step{-}by{-}step explanation and reasoning, whereas the required format for the final answer is simply the option letter (e.g., "A: (B)").} \\
            }
        }
    }
    & \\ \\

    \theutterance \stepcounter{utterance}  
    & \multicolumn{4}{p{0.6\linewidth}}{
        \cellcolor[rgb]{0.561,0.737,0.561}{
            \makecell[{{p{\linewidth}}}]{
                \texttt{\tiny{[P1$\rangle$GM]}}
                \texttt{A: (B)} \\
            }
        }
    }
    & & \\ \\

    \theutterance \stepcounter{utterance}  
    & & & \multicolumn{2}{p{0.3\linewidth}}{
        \cellcolor[rgb]{0.9,0.9,0.9}{
            \makecell[{{p{\linewidth}}}]{
                \texttt{\tiny{[GM$|$GM]}}
                \texttt{Target: (B)} \\
            }
        }
    }
    & & \\ \\

    \theutterance \stepcounter{utterance}  
    & & & \multicolumn{2}{p{0.3\linewidth}}{
        \cellcolor[rgb]{0.9,0.9,0.9}{
            \makecell[{{p{\linewidth}}}]{
                \texttt{\tiny{[GM$|$GM]}}
                \texttt{game\_result = LOSE} \\
            }
        }
    }
    & & \\ \\

\end{supertabular}
}

\end{document}
